% Context from chapters-tex/shannon-coding.tex; for mathematical reference, not compilation.
\section{Shannon Source Coding Theorem}

For a symbol distribution $p$, the expected codeword length of a code $c$ is
\eq{
	L(c) \eqdef \sum_k p_k |c_k|. 
}
We seek a prefix code minimizing $L(c)$.
 
Shannon's source coding theorem bounds the optimal expected length in terms of entropy.

\begin{thm}\label{thm-shannon-source-coding}
Assume $p_k>0$ for every symbol in the alphabet. Symbols of zero probability may be omitted.
	(i) If $c=\text{Leaves}(T)$ for some tree $T$, then 
	\eq{
		L(c) \geq H(p).
	}
	(ii) Conversely, there exists a code $c$ with $c=\text{Leaves}(T)$ such that 
	\eq{
		L(c) < H(p)+1.
	}
\end{thm} 

Kraft's inequality characterizes the possible codeword lengths of a prefix code and is the key to the proof.

\begin{lem}[Kraft inequality]\label{lem-kraft}
	(i) For a code $c$, if there exists a tree $T$ such that $c=\text{Leaves}(T)$ then
	\eql{\label{eq-kraft-ineg}
		\sum_k 2^{-|c_k|} \leq 1.
	}
	(ii) Conversely, if $(\ell_k)_k$ are nonnegative integers such that
	\eql{\label{eq-kraft-ineg-2}
		\sum_k 2^{-\ell_k} \leq 1
	}
	then there exists a code $c=\text{Leaves}(T)$ such that $|c_k|=\ell_k$.
\end{lem} 

\begin{proof}
	$\Rightarrow$ Suppose $c=\text{Leaves}(T)$ and put $m=\max_k |c_k|$. Complete $T$ to a full binary tree of depth $m$.
	 
	The subtree rooted at $c_k$ has height $m-|c_k|$ and $2^{m-|c_k|}$ leaves; see Figure~\ref{fig-kraft-trees}, panel 1. These subtrees are disjoint, and the full tree has $2^m$ leaves, so
	\eq{
		\sum_k 2^{m-|c_k|} \leq 2^m, 
	}
	which gives~\eqref{eq-kraft-ineg} after division by $2^m$.
	
	\begin{figure}[tbp]
\centering
{\small\sffamily\color{BookMuted}Panel 1\par}\smallskip
\BookDrawing{tikz/coding-kraft-necessity}
\par\medskip
{\small\sffamily\color{BookMuted}Panel 2\par}\smallskip
\BookDrawing{tikz/coding-kraft-sufficiency}
\caption{Panel 1: full binary tree obtained by completing the tree associated with the code $(c_1=0, c_2=10, c_3=110, c_4=111)$.
	Panel 2: subtrees with the prescribed codeword lengths packed from left to right.}
\label{fig-kraft-trees}
\end{figure}

	
	$\Leftarrow$ Conversely, assume~\eqref{eq-kraft-ineg-2} holds and order the lengths so that $\ell_1 \leq \ldots \leq \ell_K =: m$.
	We construct full subtrees of height $m-\ell_k$, each with $2^{m-\ell_k}$ leaves, and place their leaf blocks consecutively from left to right.
	 
	The inequality $\sum_k 2^{m-\ell_k} \leq 2^m$ ensures that all blocks fit among the $2^m$ leaves.
	 
	Because the lengths $\ell_k$ are nondecreasing, the block sizes are nonincreasing powers of two. Each block therefore begins at a multiple of its own size and is the leaf set of a subtree; see Figure~\ref{fig-kraft-trees}, panel 2. The roots of these disjoint subtrees form the required prefix code.
\end{proof}



\begin{proof}[Proof of Shannon's theorem]
(i) Put $\ell_k=|c_k|$, $Z=\sum_k2^{-\ell_k}\leq1$, and $q_k=2^{-\ell_k}/Z$. The nonnegativity of relative entropy gives
\eql{\label{eq-proof-shannon}
 L(c)-H(p)=\mathrm{KL}(p\mid q)-\log_2 Z\geq0.
}
Equality requires $Z=1$ and $q=p$, identifying the ideal real-valued code lengths $-\log_2p_k$.

(ii) Choose $\ell_k=\lceil-\log_2p_k\rceil$. Then $\sum_k2^{-\ell_k}\leq\sum_kp_k=1$, so Lemma~\ref{lem-kraft} constructs a prefix code with these lengths. Moreover,
\eq{
 L(c)=\sum_kp_k\lceil-\log_2p_k\rceil
 <\sum_kp_k(-\log_2p_k+1)=H(p)+1.
}
For a one-symbol alphabet, the empty word has length zero and the same conclusions hold.
\end{proof}

This proof is constructive: rounding the ideal lengths and using Kraft's lemma yields a Shannon code. It need not minimize $L(c)$. Huffman's greedy algorithm constructs an optimal binary prefix code by repeatedly merging the two least probable symbols. With a priority queue, its cost is $O(K\log K)$. Figure~\ref{fig-huff-code} illustrates the construction.


\myfigure{
\includegraphics[width=.8\linewidth]{1-shannon/huff-code}
}{Successive merges in Huffman's coding algorithm.}{fig-huff-code}

\texttt{Associated code: coding/test\_text.m}

                                                            
\secti